% Part: first-order-logic
% Chapter: introduction
% Section: satisfaction

\documentclass[../../../include/open-logic-section]{subfiles}

\begin{document}

\olfileid{fol}{int}{sat}

\olsection{صدق}

کله چې پوه شو !!{formula}s څه دي، سمدستي د لومړۍ درجې منطق معنٰی
پوهنې ته وړاندې تللے شو: دلته بنسټيز تعريف د !!a{structure} دے۔
زموږ د ساده ژبې دپاره !!a{structure}~$\Struct M$ يوازې درې برخې
لري: يو ناتش سټ~$\Domain{M}$ چې \emph{!!{domain}} بلل کېږي، هغه
څيز چې~$\Obj a$ ئې په~$\Struct M$ کښې ټاکي، او هغه څيزونه چې
$\Obj P$ ورباندې په~$\Struct M$ کښې رښتونے دے۔ د~$\Obj a$ ټاکلے
څيز په~$\Assign{\Obj a}{M}$، او د هغو څيزونو سټ چې~$\Obj P$
ورباندې رښتونے دے په~$\Assign{\Obj P}{M}$ ښودل کېږي۔
!!^a{structure}~$\Struct{M}$ همدا درې څيزونه لري: $\Domain{M}$،
$\Assign{\Obj a}{M} \in \Domain{M}$ او
$\Assign{\Obj P}{M} \subseteq \Domain{M}$۔ عمومي حالت به لا پېچلے
وي، ځکه ډېر !!{predicate}s او !!{constant}s به وي،
پريديکاتونه له يوه څخه زيات ځايونه لرلے شي، او !!{function}s به
هم وي۔ \textbf{د ژباړې د نسخې سمون (OLFOL-002):} سرچينې دلته په
تېروتنه کښې ويلي چې !!{constant}s له يوه څخه زيات ځايونه لرلے شي؛ د
څوځاييزوالي خاصيت د پريديکاتونو دے، لکه د همدې څپرکي دوه‌ځاييزه
اړيکه۔

دا د هغو !!{formula}s د صدق د تعريف دپاره بس ده چې !!{variable}s
نۀ لري۔ مفکوره دا ده چې داسې استقرايي تعريف ورکړو چې د
!!{formula}s زموږ د تعريف طريقه منعکس کړي۔ ټاکو چې يو اټومي فارمول
کله په~$\Struct{M}$ کښې صدق کوي؛ بيا مثلاً دا ټاکو چې $\lnot !A$
کله په~$\Struct{M}$ کښې صدق کوي، د دې پر بنسټ چې~$!A$ په~$\Struct{M}$ کښې صدق
کوي که نۀ۔ مثلاً داسې تعريفولے شو:
\begin{enumerate}
  \item $\Atom{\Obj P}{\Obj a}$ په~$\Struct{M}$ کښې هله صدق کوي
  چې $\Assign{\Obj a}{M} \in \Assign{\Obj P}{M}$ وي۔
  \item $\lnot !A$ په~$\Struct{M}$ کښې هله صدق کوي چې $!A$ په
  ~ $\Struct{M}$ کښې صدق نۀ کوي۔
  \item $(!A \land !B)$ په~$\Struct{M}$ کښې هله صدق کوي چې $!A$
  په~$\Struct{M}$ کښې صدق وکړي، او $!B$ هم په~$\Struct{M}$ کښې
  صدق وکړي۔
\end{enumerate}
فرض کړئ $\Domain{M} = \{0, 1, 2\}$، $\Assign{\Obj a}{M} = 1$، او
$\Assign{\Obj P}{M} = \{1, 2\}$ دي۔ دا تعريف راته وايي چې
$\Atom{\Obj P}{\Obj a}$ په~$\Struct{M}$ کښې صدق کوي (ځکه
$\Assign{\Obj a}{M} =  1 \in \{1,2\} = \Assign{\Obj P}{M}$)۔
ورپسې راته وايي چې $\lnot \Atom{\Obj P}{\Obj a}$ په~$\Struct{M}$
کښې صدق نۀ کوي؛ نو په خپل وار $\lnot\lnot \Atom{\Obj P}{\Obj a}$
صدق کوي او $(\lnot \Atom{\Obj P}{\Obj a} \land \Atom{\Obj P}{\Obj a})$
صدق نۀ کوي، او همداسې مخکښې۔

ستونزه هغه وخت پيدا کېږي چې د کميت ټاکونکو دپاره تعريف ورکول
غواړو: زړه مو غواړي داسې ووايو چې
``$\lexists[\Obj v_0][\Atom{\Obj P}{\Obj v_0}]$ هله صدق کوي چې
$\Atom{\Obj P}{\Obj v_0}$ صدق وکړي''۔ خو !!{structure}~$\Struct{M}$
راته نۀ وايي چې له !!{variable}s سره څه وکړو۔ په اصل کښې غواړو
ووايو چې $\Atom{\Obj P}{\Obj v_0}$ \emph{د~$\Obj v_0$ د کوم قيمت
دپاره} صدق کوي۔ د دې د دقيق کولو دپاره داسې لارې ته اړتيا لرو چې
د~$\Domain{M}$ !!{element}s يوازې~$\Obj a$ ته نۀ، بلکې~$\Obj v_0$
ته هم وګومارو۔ د همدې دپاره د !!{variable} \emph{ګومارنې} ورپېژنو۔
د !!^a{variable} ګومارنه يوازې داسې تابع~$s$ ده چې !!{variable}s
د~$\Domain{M}$ !!{element}s ته وړي (زموږ په بېلګه کښې له~$0$، $1$
يا~$2$ څخه يوه ته)۔ \textbf{د ژباړې د نسخې سمون (OLFOL-003):}
سرچينې دلته يو، دوه او درې ليکلي وو، خو پورته ټاکلې د تعريف ساحه
صفر، يو او دوه لري؛ نو لړ د هماغې ساحې غړو ته سم شو۔ ځکه مخکښې نۀ
پوهېږو چې په !!a{formula} کښې به کوم !!{variable}s راشي، دا نۀ شو
محدودولے چې~$s$ کومو !!{variable}s ته قيمتونه وټاکي۔ ساده حل دا
دے چې وغواړو~$s$ \emph{ټولو} !!{variable}s~$\Obj v_0$،
$\Obj v_1$، \dots\@ ته قيمتونه وګوماري۔ موږ به يوازې د اړينو هغو
قيمتونه کاروو۔

د دې پر ځاے چې د !!{formula}s صدق يوازې د !!a{structure} په نسبت
تعريف کړو، دا به د !!a{structure}~$\Struct{M}$ \emph{او} د
!!a{variable} ګومارنې~$s$ دواړو په نسبت تعريف کړو، او په لنډ ډول به
$\Sat{M}{!A}[s]$ ليکو۔ اوس به زموږ تعريف د هغو اټومي !!{formula}s
دپاره يو اضافي شرط هم ولري چې !!{variable}s لري:
\begin{enumerate}
  \item $\Sat{M}{\Atom{\Obj P}{\Obj a}}[s]$ هله او يوازې هله چې
  $\Assign{\Obj a}{M} \in \Assign{\Obj P}{M}$ وي۔
  \item $\Sat{M}{\Atom{\Obj P}{\Obj v_i}}[s]$ هله او يوازې هله چې
  $s(\Obj v_i) \in \Assign{\Obj P}{M}$ وي۔
  \item $\Sat{M}{\lnot !A}[s]$ هله او يوازې هله چې
  $\Sat{M}{!A}[s]$ نۀ وي۔
  \item $\Sat{M}{(!A \land !B)}[s]$ هله او يوازې هله چې
  $\Sat{M}{!A}[s]$ او $\Sat{M}{!B}[s]$ دواړه وي۔
\end{enumerate}
ښه، په دې سره يوه ستونزه هواره شوه: اوس ويلے شو چې~$\Struct{M}$ د
$s(\Obj v_0)$ قيمت دپاره $\Atom{\Obj P}{\Obj v_0}$ کله پوره کوي۔
د $\lexists[\Obj v_0][\Atom{\Obj P}{\Obj v_0}]$ د تعريف د سمولو
دپاره يو بل کار پکار دے: غواړو
$\Sat{M}{\lexists[\Obj v_0][\Atom{\Obj P}{\Obj v_0}]}[s]$ هله او
يوازې هله وي چې د~$\Obj v_0$ د قيمت د ګومارلو د \emph{کومې} لارې
$s'$ دپاره $\Sat{M}{\Atom{\Obj P}{\Obj v_0}}[s']$ وي۔ خو~$\Obj v_0$
ته ګومارلے قيمت لازماً هغه قيمت نۀ دے چې~$s(\Obj v_0)$ ئې ټاکي۔
د دې دپاره يوه نښه ورپېژنو: که $m \in \Domain{M}$ وي، نو
$\Subst{s}{m}{\Obj v_0}$ هغه ګومارنه ده چې (له~$\Obj v_0$ پرته د
ټولو !!{variable}s دپاره) کټ مټ د~$s$ په شان وي، خو~$\Obj v_0$ ته
$m$ ګوماري۔ اوس زموږ تعريف داسې کېدے شي:
\begin{enumerate}\setcounter{enumi}{4}
  \item $\Sat{M}{\lexists[\Obj v_i][!A]}[s]$ هله او يوازې هله چې
  د کوم~$m \in \Domain{M}$ دپاره
  $\Sat{M}{!A}[\Subst{s}{m}{\Obj v_i}]$ وي۔
\end{enumerate}
ايا دا کار کوي؟ فرض کړئ د ټولو~$i \in \Nat$ دپاره
$s(\Obj v_i) = 0$ ټاکو۔
$\Sat{M}{\lexists[\Obj v_0][\Atom{\Obj P}{\Obj v_0}]}[s]$ هله او
يوازې هله چې داسې~$m \in \Domain{M}$ وي چې
$\Sat{M}{\Atom{\Obj P}{\Obj v_0}}[\Subst{s}{m}{\Obj v_0}]$ وي۔ او
داسې غړے شته: $m = 1$ يا $m = 2$ ټاکلے شو۔ پام وکړئ دا حتٰی هغه
وخت رښتيا دي چې~$s$ پخپله~$\Obj v_0$ ته ګومارلے قيمت
$s(\Obj v_0)$---دلته~$0$---کار نۀ کوي۔
$\Sat{M}{\Atom{\Obj P}{\Obj v_0}}[\Subst{s}{1}{\Obj v_0}]$ لرو، خو
$\Sat{M}{\Atom{\Obj P}{\Obj v_0}}[s]$ نۀ لرو۔

که دا مغشوش او دروند ښکاري، همداسې دے۔ خو دا زياته پېچلتيا د ټولو
!!{formula}s د صدق د دقيق، استقرايي تعريف دپاره اړينه ده، او د
معنايي مفاهيمو د دقيق تعريف دپاره ورته څه پکار دي۔ نورې لارې هم
شته، خو ټولې په همدې کچه ښکلې (يا ناشاعرانه) دي۔

\end{document}
